% !PRINTED 561
\weberpage{561}
%<<CONT
dinaten eines Punktes $\xi$ in einem Raume von $n$ Dimensionen
$R_\xi$ und nennen den Ausdruck

\weq{7}{\varrho = \sqrt{(\xi_1' - \xi_1'')^2 + (\xi_2' - \xi_2'')^2 + \cdots + (\xi_n' - \xi_n'')^2},}

die \gesperrt{Entfernung} der beiden Punkte $\xi'$, $\xi''$.

Wenn wir in (3) die Variablen $x_i$ alle ganzzahligen Werthe
durchlaufen lassen, so geben uns die zugehörigen Werthe der $\xi_i$
ein Punktsystem im Raume $R_\xi$. Der Punkt, der den Werthen $x_i =
0$ entspricht, ist der Coordinaten-Anfangspunkt oder der \gesperrt{Nullpunkt}.
Die Gesammtheit dieser Punkte bleibt ungeändert, wenn wir $x_i$
durch $x_i + x_i^{(0)}$ ersetzen, wenn $x_1^{(0)}, x_2^{(0)}, \ldots,
x_n^{(0)}$ ein festes System ganzer Zahlen bedeutet. Wir können diese
Thatsache auch so ausdrücken, dass das Punktsystem der $\xi_i$ mit
sich selbst zur Deckung kommt, wenn es eine Parallelverschiebung
erfährt, bei dem der Nullpunkt an den Ort eines beliebig gegebenen
anderen Punktes des Systems gelangt; oder auch: Die Punkte des
Systems sind zu jedem seiner Punkte in derselben Weise gelagert, wie
zu jedem anderen.

Wir wollen dies System ein \gesperrt{Punktgitter} nennen und seine
einzelnen Punkte die \gesperrt{Gitterpunkte}.

Nach (2) und (7) ist der Werth von $f$ das Quadrat der Entfernung des
Punktes $\xi$ vom Nullpunkte und folglich ist $\sqrt{M}$ die
\gesperrt{kleinste Entfernung zweier Gitterpunkte}, die in dem ganzen System
vorkommt.

Ausser den Entfernungen in dem Raume $R_n$ haben wir auch noch
Volumina zu betrachten, d.\,h. die Werthe des $n$-fachen Integrals

\weq{8}{V_\xi = \int\!\!\int \ldots \int d\xi_1\, d\xi_2 \ldots d\xi_n,}

worin die Variablen $\xi_1, \xi_2, \ldots, \xi_n$ irgend ein
endliches Gebiet $G_\xi$ erfüllen, was etwa durch eine oder mehrere
Ungleichungen bestimmt ist. So ist das durch die Bedingung

\weq{9}{\xi_1^2 + \xi_2^2 + \cdots + \xi_n^2 \leqq 1 \text{ oder auch } \Sigma(\xi_i - \xi_i^0)^2 \leqq 1}

bei beliebigen Werthen $\xi_i^0$ begrenzte Gebiet eine
\gesperrt{$n$-dimensionale Kugel} vom Radius 1, deren Volumen wir mit $K_n$
bezeichnen wollen. Aus jedem Gebiete $G_\xi$ können wir ein
\gesperrt{ähnliches Gebiet} $G_\xi^{(t)}$ ableiten, indem wir die Coordinaten
alle im Verhältniss $t^{\frac{1}{n}} : 1$ vergrössern. Ist $V^{(t)}$
das Volumen des neuen Gebietes, so ist nach (8)

\weq{10}{V^{(t)} = t\,V.}
